\section{Background and Related Works}
\label{sec:related_works}
Our work lies at the intersection of research directions addressing the inherent trilemma of FL: statistical heterogeneity, system heterogeneity, and communication bottlenecks. The first category, standard FL, aims to collaboratively train global models without centralizing raw data. Fundamental algorithms like FedAvg~\cite{pmlr-v54-mcmahan17a} often suffer from severe ``client drift'' due to non-IID data distributions. To mitigate this, traditional solutions rely on regularization mechanisms, such as proximal terms in FedProx~\cite{dai2023tacklingFedProx}, or momentum-based optimization~\cite{cheng2024momentumbenefits, zaccone2025communication}. To address variable local work across heterogeneous clients, FedNova~\cite{10.5555/3495724.3496362} normalizes client updates before aggregation. However, a critical limitation of these frameworks is that they inherently assume clients possess the computational capacity to train and evaluate full models locally. This assumption exacerbates system heterogeneity and excludes resource-constrained edge devices from participating.

To alleviate these client-side computational constraints, SFL partitions the model and offloads heavy upper layers to the server. SplitFed~\cite{thapa2022splitfed} was introduced to parallelize client updates, resolving the sequential training bottlenecks of early split learning. This literature has expanded to address computational diversity and scale through HeteroSFL paradigms that build on sub-model deployments such as HeteroFL~\cite{diao2021heterofl} and Hierarchical SFL (HSFL)~\cite{rana2022hierarchical,9148862,9488756}, which introduce edge servers for intermediate synchronization. Despite these topological advancements, these frameworks suffer from limitations. They require the continuous transmission of high-dimensional activations during the forward pass and the subsequent return of gradients during the backward pass. This bidirectional dependency introduces synchronization overhead, exacerbates communication bottlenecks, and creates new computational bottlenecks on the server~\cite{10980018}.

As an alternative to raw gradient and activation communication, the third category explores generative and prototype-based FL. Methods such as Federated Knowledge Distillation (FedDF~\cite{lin2020ensemble}) and data-free Generative FL (FedGen~\cite{venkateswaran2023fedgen}) address both communication overhead and non-IID challenges but typically require training computationally expensive generators, such as GANs, at the server. Conversely, prototype-based FL, such as FedProto~\cite{tan2022fedproto} or FOUL \cite{nguyen2026computation}, exchanges lightweight class prototypes instead of gradients or heavy activations. A major limitation of prototype methods, however, is that they primarily utilize these prototypes as client-side regularization terms to align local feature spaces. Because they are not used to synthesize actual training data, these methods can still suffer from prototype drift under extreme heterogeneity~\cite{bensiah2025feddpa, zhang2020personalized}.

\textcolor{black}{H-SFP complements these resource-efficient and decoupled learning approaches. PriSM~\cite{niu2023overcoming} reduces client computation by sampling principal kernels to construct lightweight sub-models, while Delta~\cite{10655220} studies asymmetric private/public computation flows with explicit privacy mechanisms. In contrast, H-SFP focuses on the HSFL setting, where clients, edge servers, and the cloud train different model segments. Its main design choice is a statistic-based cross-tier interface: upper tiers receive class-wise feature moments rather than sample-wise activations, gradients, full model updates, or learned generators. This makes H-SFP particularly suitable for hierarchical deployments.}

Existing frameworks preserve a shared end-to-end optimization graph across tiers. Even when activations or prototypes are transmitted rather than full models, the upper and lower tiers remain coupled via cross-tier gradient propagation or parameter synchronization. This architectural dependency forces all components to participate in a single distributed optimization process, making communication and statistical drift intrinsically entangled.
In contrast, H-SFP departs from this paradigm by severing the optimization graph across tiers. \textcolor{black}{Rather than coordinating updates through sample-wise gradients or activations, H-SFP performs tier-wise optimization over synthesized feature distributions derived from aggregated statistical moments. This design provides a lightweight communication interface between client, edge, and cloud tiers while preserving the practical advantage of split-model execution on resource-constrained clients.}
A detailed comparison between approaches is shown in the Appendix~A.
